Markov Chain Monte Carlo (MCMC)

What is a (Markov) kernel? What is a probability kernel?

Suppose that \( (S, \mathcal{S}) \) and \( (T , \mathcal{T}) \) are measurable spaces. Then
\( µ : S \times \mathcal{T} \to \mathbb{R}_+ \cup \{\infty\}\)

is called a (Markov) kernel from \( S \) to \( T \) if
1. \( \forall B \in \mathcal{T}, s \mapsto \mu_s(B) = \mu(s, B) \) is \( \mathcal{S} \)-measurable.
2. \( \forall s \in S, s \mapsto \mu_s \) is a measure on \( \mathcal{T} \)
The special case of a probability kernel, refers to the situation \( \mu : S \times \mathcal{T} \to [0, 1] \) 
such that \( \mu_s \) is a probability measure for all \( s \in S \).


Define the notion of random elements.

Let \((\Omega, \mathcal{F}, P)\) be a probability space, and \( (E, \mathcal{E}) \) a measurable space. A random element with values in \( E \) is a function \( X : \Omega \to E \) which is 
\((\mathcal{F}, \mathcal{E})\)-measurable. That is, a function X such that for any \( B \in \mathcal{E} \), the preimage of \( B \) lies in \( \mathcal{F} \).

Sometimes random elements with values in \( E \) are called \( E \)-valued random variables.

Note if \( (E, \mathcal{E}) = (\mathbb{R} , \mathcal{B}(\mathbb{R}))\), where \( \mathcal{B}(\mathbb{R}) \) is the Borel \( \sigma \)-algebra, 
then the definition of random element is the classical definition of random variable.

The power lies in the abstraction, since \( E \) could also be a space of functions (e.g. \(E = C[0, 1] \) with norm \( \|f\| = \sup_{x \in [0, 1]} |f(x)| \) and the corresponding Borel \( \sigma \)-algebra, which is typically the setup for continuous time stochastic processes.

Define the notion of a regular conditional distribution. What is a sufficient
condition for its existence?

Let \( (\Omega, \mathcal{A}, \mathbb{P}) \) be a probability space and \( (S, \mathcal{S}) \), \( (T , \mathcal{T}) \) measurable spaces.
Given random elements \( \xi : \Omega \to T \), \( \nu : \Omega \to S\)
it may be possible to define a conditional distribution of \( \xi \) given \( \sigma(\eta) \), that is
\( \mathbb{P}[\xi \in A|\eta], A \in \mathcal{T} \) and in particular we would like to have a probability kernel
with \( \mu(\eta, A) = \mathbb{P}[\xi \in A|\eta] \).

Its existense is guaranteed if \( T \) is Polish (completely metrizable separable space).


Define a Markov process. What is the extended Markov property?

The process \( X : (\Omega, I) \to E\) has the Markov property, if
for all \( A \in \mathcal{B}(E) \) and all \( s, t \in I \) with \( s \leq t \)
1. \( X \) is \( \mathbb{F} \)-adapted.
2. \( \mathbb{P}[X_t \in A | \mathcal{F}_s] = \mathbb{P}[X_t \in A | X_s]\)

In the lecture (2) is also expressed as \( \mathcal{F_t} \) being conditionally independent of \( X_s \) given \( X_t \).
The idea of requirement (2) is that \( \mathcal{F}_s \) contains all events that have potentially
lead to the state \( X_t \), while \( |X_s = \sigma(X_s) \) refers only to the specific event, which causes the value at \( X_s \)
(so all \(X_k, k < s \) do not matter).

The strong/extended Markov property asserts that given a stopping time \( \tau \), \( \mathbb{P}[X_{\tau + t} \in A | \mathcal{F}_\tau] = \mathbb{P}[X_{\tau + t} \in A | X_\tau] \),
equivalently \( \mathcal{F_\tau} \) is conditionally independent of \( X_{\tau + t} \) given \( X_{\tau} \) for all \( t \in I \).

What are the Chapman-Kolmogorov relations? State Kolmogorov’s existence theorem.

A family of finite-dimensional distributions is defined as:
Let \( X = \{X_t\}_{t \in T} \) a Markov process with one-dimensional distributions \( \nu_t \) and transition kernels \( \mu_{s,t} \),
\( s, t \in T \), \( s \leq t \). Then for all \( t_0 \leq \dots \leq t_n \in T \) :
\( \mathcal{L}(X_{t_0} , X_{t_1}, \dots, X_{t_n}) = \nu_{t_0} \otimes \mu_{t_0,t_1} \otimes \dots \otimes \mu_{t_{n-1}, t_n} \)
\( \mathbb{P}[(X_{t_1}, \dots, X_{t_n}) \in \cdot | F_{t_0}] = (\ \mu_{t_0,t_1} \otimes \dots \otimes \mu_{t_{n-1}, t_n}(X_{t_0}, \cdot) \)

Then Kolmogorov’s existence theorem states:
Suppose an initial distribution \( \nu_{t_0} \) and transition kernels \( \mu_{s,t} \) are given satisfying the Chapman-Kolmogorov
relations. Then there exists a Markov process with these specifications.

The Chapman-Kolmogorov relations are:
\( \mu_{s,u} = \mu_{s, t}\mu{t,u}, s \leq t \leq u \)


What is a time-homogeneous Markov process? State the Chapman-Kolmogorov
relations in this case, also called the semigroup property.

Let \( T = \mathbb{R}_{+} \) or \( T = \mathbb{Z}_{+} \).
A time-homogeneous Markov process possesses transition kernels \( \mu_{s,t} \) of the form \( \mu_{t-s} \) , s.t.
\( \mathbb{P}[X_t \in B|\mathcal{F}_s] = \mu_{t-s}(X_s , B) \) a.s., \( B \in \mathcal{S}\), \( s \leq t \).

Then the finite-dimensional distributions simplify to 

\( \mathcal{L}(X_{t_0} , X_{t_1}, \dots, X_{t_n}) = \nu \mu_{t_0} \otimes \mu_{t_1 - t_0} \dots \otimes \mu_{t_n - t_{n-1}} \)
\( \mathbb{P}[(X_{t_1}, \dots, X_{t_n}) \in \cdot | F_{t_0}] = (\ \mu_{t_1 - t_0} \otimes \dots \otimes \mu_{t_n - t_{n-1}}(X_{t_0}, \cdot) \)

and the Chapman-Kolmogorov relations are:
\( \mu_{s+t} = \mu_{s}\mu{t}, s, t \in T \)
which is also referred to as semigroup property.


Explain how the law of a Markov process may be derived:

The full question reads:
A Markov process is a stochastic process (family of random elements) on a filtered
probability space with specific properties (as stated). Alternatively, the process
can be interpreted as a random element in a suitable product space. How does
the construction work? What is the law of the process? What can you say about
10uniqueness? How can the law of the Markov process for a fixed transition semigroup
be reconstructed for any initial distribution if its law is known for deterministic
starting values (i.e., Dirac measures)?

Let \( S^{T} \) be the space of all maps \( f \) from \( T \) to \( S \) with the \( \sigma \)-algebra
generated by the projection \( \pi_t(f) = f(t) \), that is \( \mathcal{S}^{T} = \sigma(\{\pi_t : S^{T} \to S, t \in T\}) \)
It can be proven that for a fixed initial distribution \( \nu \) on \( S \) and semigroup of transition kernels \( \mu_{s,t} \)
\( X_{\nu} : \Omega \to S^{T}, \omega \mapsto (X_{\nu,t}(\omega)) \)
is measurable with respect to \( \mathcal{S}^{T} \), so its a random element in \( S^{T}, \mathcal{S}^T \).
The image measure of \( X_\nu \) is the distribution of the Markov process \( \mathbb{P}_\nu = \mathcal{L}(X_\nu) \),
which is unique, even though multiple \( X_\nu \) may have this law.

In particular this is a mixture \( \mathbb{P}_\nu = \int_{S} \mathbb{P}_{\delta_x}(A)\nu(dx)\).

What is the canonical construction of a stochastic process?

Recall the design we leveraged for the law of a Markov process then
the canonical process on \( (S^T , \mathcal{S}^T) \)
is defined by setting \( X = \text{id}, \omega \mapsto \omega \)
(so \( \omega \) is a function from \( T \) to \( S \) with the projection 
\( \pi_t(\omega) = \omega(t) \)) 
with the natural/canonical filtration \( \mathcal{F}_t = \sigma(\{X_s\}_{s \leq t}) \) then
\( X : (S^{T}, \mathcal{F}_t, \mathbb{P}_\nu) \to (S^T, \mathcal{S}^T) \), where  
\( \mathbb{P}_\nu \) is the law determined uniquely by \( \nu \), \( \mu_{s,t} \),
is the so called canonical construction.


What is a statistical model? What is the difference between the classical approach
and Bayesian statistics? State Bayes’ formula. Provide application examples.

A statistical model \( (\Chi , \mathcal{F}, \mathbb{P}_\vartheta : \vartheta \in \Theta) \) consist of
- The space of observable data \( \Chi \) (similar to the usual sample space \( \Omega \))
- The \( \sigma \)-algbebra of events \( F \)
- The admissable models \( \{\mathbb{P}_\vartheta | \vartheta \in \Theta \} \) 
- The parameter space \( \Theta \)


Two common approaches to determine \( \{\mathbb{P}_\vartheta | \vartheta \in \Theta \} \)
are the classical/frequentist approach and the Bayesian one.

The classical approach usually assumes that there is only one correct \( \vartheta \),
which can be determined by a large number of observations.

The Bayesian instead assumes a probability distribution \( \pi \sim \vartheta \) on \( \Theta \),
called the prior distribution.
It is also derived by a large number of observations and some assumptions about the distribution.

In the Bayesian approach we do not have direct access to \( \pi \), but only the posterior distribution
\( \pi(\vartheta|x) = \frac{\pi(\vartheta)p(x|\vartheta)}{p(x)} = (\int_{\Theta} \pi(\vartheta)p(x|\vartheta) d\vartheta)^{-1} \pi(\vartheta)p(x|\vartheta) \)
(this is an application of Bayes rule)
where \( p(x) \) is the probability density function of the observed data \( X \) and \( p(x|\vartheta) \) the conditional density on \( \vartheta \)
which receives the special name likelihood in this context.
\( \mathbb{P}[x, \vartheta] = \pi(\vartheta)p(x|\vartheta) \) is the joint distribution.

The most famous real-world application is probably in medical studies:
Suppose we developed an early detection test for lung cancer and want
to show its efficiancy. 
Let \( \Chi = \{B, \overline{B}\} \) with \( B \) representing a positive test
and \( \overline{B} \) an inconclusive/negative test result.

Let \( \Theta = \{A, \overline{A}\} \), where \( A \)
stands for lung cancer and \( \overline{A} \) for no disease.

We are then interested in \( \pi(\vartheta = A|x = B) \) that is the probability
of a patient having cancer, given that they had a positive test.

Note that we have direct access to all the other information \( p(x = B) \)
is simply the observed frequency of positive/negative results disregarding if the subject actually had cancer.
\( \pi(\vartheta = A) \) the probability of the general population to have lung cancer 
(or the ratios of our control and test groups in the study),  
and \( p(x|\vartheta = A) \) is the share of true positives, false positives in our study.

Actually \( p(x = B) \) may not be straightforward to determine
and that's where Markov Chain Monte Carlo steps in.


Justify that for the simulation of the posterior using MCMC we only need to
understand the numerator in Bayes’ formula. Or more generally explain the big
picture behind MCMC.

This also includes the answer to:
Consider the inverse problem of independent normal data with known variance and
uncertain mean. Suppose that the prior for the mean is also normal. Compute the
posterior using the “proportionality trick”.


Note that by Bayes formula \( \pi(\vartheta|x) \propto \pi(\vartheta)p(x|\vartheta) = f_{\vartheta, \pi} \)
so the posterior is proportional to the joint distribution, where
\( \propto \) refers to equality up to multiplication by a constant. 

Afterall 
we have access to the observations \( x \) and so the integral \( p(x) = \int_{\Theta} \pi(\vartheta)p(x|\vartheta) d\vartheta \)
directly evaluates to a constant and is not dependent on \( \vartheta \) during evaluation.

So in general MCMC tries to solve the problem, that we want to sample from a distribution \( f(z) = \frac{1}{Z_f}\tilde{f}(z) \)
where \( \tilde{f}(z) \) is easy to compute, but \( \frac{1}{Z_f} \) hard or near impossible.

There is the special situation of conjugate priors, where the posterior 
has the same distribution as the prior, so \( \pi(\vartheta|x;\alpha_1) ~ \pi(\vartheta;\alpha_0) \)
but potentially with a different set of so called hyperparameters \( \alpha_0(x) \to \alpha_1(x) \) which may depend on \( x \).
We say that the likelihood function \( p(x|\vartheta \) is conjugate prior to \( \pi(\vartheta) \) and we have conjugate famility.
E.g. \( \pi(\vartheta) ~ N(\mu,\sigma^2) \) then \( (\mu,\sigma^{2}) \) are the hyperparameters \( \alpha_0 \)

More concretely:
Suppose \( \vartheta \sim \mathcal{N}(\mu_0, \gamma_2) \) and \( p(X_i | \vartheta) = \mathcal{N}(\vartheta, \sigma_2) \)
for \( i = 1, \dots, N \) with known \( \sigma_2 \) and of course unknown \( \vartheta \). In this case we
\( α_0 = (µ_0, γ^2) \). If \( X = (X_1, \dots, X_N ) \).
We then have
\( p(θ | x) \propto p(x | θ)\pi(θ; α0) \propto \exp{-\frac{1}{2\gamma_0^2}(\vartheta - \mu_)^2}\Pi_{i=1}^{N} \exp(-\frac{1}{2\sigma^2}(x_i - \vartheta)^2 \)
\( \propto \exp(-\frac{1}{2\gamma_1^{2}}(\vartheta - \mu_1)^2 \)
where \( \gamma_1^2 = (\gamma_0^{-2} + N\sigma^{-2})^{-1} \) and \( \mu_1 = \gamma_1^{2}(\mu_0 \gamma_0^{-2} + \sum_{i=1}^N x_i \sigma^{-2} \)
We recognize \( p(\vartheta|x) \sim N(\mu_1, \gamma_2\)

Note that if we let \( N \to \infty \) then \( \gamma_1 \to 0 \) and \( (\vartheta - \mu_1)^{2} \to \vartheta \).

How can particle systems (static situation) be modeled? Define the Ising model.
Why are direct computations based on its distribution typically infeasible – im-
plying that the application of MCMC is important?

We consider a discrete set of sites \( S \).
The sites are in certain states, e.g., in \( \{-1, 1\} \).
This leads to configurations \( x = (x_i)_{i \in S} \) with \( x_i \in \{-1, 1\}\).
The configurations are random, i.e., there exists a probability measure on \(\{-1, 1\}^{S}\).
A canonical model can be constructed as follows:

If \( X \) is the identity on \(\{-1, 1\}^{S}\) and \( X_i : \{-1, 1\}^{S}  \to \{-1, 1\} \) is the projection on coordinate \( i \), then
\( X = (X_i)_{i \in S} \) is a random element that canonically models the random configurations.

E.g.
Ferromagnets:
site = atom, \( \{-1, 1\} \) possible spin directions
Image analysis:
\( S = \{1, 2, \dots, N\} \times \{1, 2, \dots, M\} \) set of pixels, \( {−1, +1} \) white or black.

Distributions \( \Pi^{*} \) on a configuration space can be specified in different ways.
A typical approach specifies local properties like attraction or repulsion
This requires notions of neighborhood:
As an example, consider a graph \( G = (S, E) \) with vertices \( S \) and edges \( \Eta \), i.e.,
\( i,i' \in S \) neighbors \( \iff \) \( (i, i') \in \Eta \)
Special case: \( S \in \mathbb{Z}^{2} \) with
\( (i, i') \in \Eta \iff i = (u, v), i' = (u', v'), |u' - u| + |v ′ − v | = 1\)

A well known case the Ising model

Ising model

\(  \Pi^{*}(x) \propto \pi(x) = \exp \left\{ -\beta \sum_{\{i,i'\} \in E} 1_{\{x_i \neq x_{i'}\}} \right\} \)

with inverse temperature \( \beta > 0 \)

* \( \pi(x) \) is smaller when more neighboring sites have opposite signs (local attraction).
* A lower temperature corresponding to a larger \( \beta \) makes all 1's or all -1's more likely.}

A direct computation based on \( \Pi^{*} \) is typically infeasible, since the normalizing constant cannot be computed.
This can be illustrated for a relatively small graph with \(S = \{1,2,...,8\}^2\). 

The normalizing constant involves $2^{64}$ summands:
\( C = \sum_{x \in \{-1,1\}^8} \exp \left\{ -\beta \sum_{i,j \in E} 1\{x_i \neq x_j\} \right\} \)
Illustrating the importance of MCMC.

MCMC can also be used for counting complicated combinatorial objects. Discuss
simple examples of this type.

An application of MCMC becomes also important, if $E$ is finite but very large and constructed from a complicated structure. 
What is the cardinality of \(E\)?
How can we sample from a uniform distribution on \(E\)?

Examples
1. Hard-core-model

The set \(E \subset \{-1,1\}^{\{1,2,...,N\} \times \{1,2,...,M\}}\), and \(E\) contains only those configurations such that no two neighboring points are occupied. 

2. Square ice model

The vertices are given by \(\{1,2,...,N\} \times \{1,2,...,M\}\), the edges by 
\[E = \{ii' : i = (u,v), i'=(u',v') \in V, |u'-u| + |v'-v| = 1\}.\]
Let \(S \subset P(E)\) such that \(\forall A \in S \ \forall ii' \in E: ii' \in A \vee ii' \in A\).
Consider now \(E \subset S\) such that the non-boundary vertices each possess two ingoing and two outgoing arrows (at boundary: at most).

3. Coloring

\(E\) is the set of possible allocations of \(q\) colors to boxes such that not two squares with a common boundary have the same color. 

Define a hidden Markov model. Describe applications of MCMC in this context.

A Hidden Markov model is a bivariate discrete stochastic process $(X_k, Y_k)_{k \geq 0}^\top$ such that 
    * \( X_k \) is a Markov process
    * \(Y_k\) is a sequence of random variables 
        * that are conditionally independent given \((X_k)_{k \geq 0}\), 
        * and the conditional distribution \(\mathbb{P}(Y_j \in \cdot | (X_k)_{k \geq 0}) = \mathbb{P}(Y_j \in \cdot | X_j)\) depends only on \(X_j\), \(j \geq 0\)

Hidden Markov models admit a functional representation as a general state-space model, i.e., 
\[
X_{k+1} &= a(X_k, U_k) \\
Y_k &= b(X_k, V_k)
\]

with measurable functions \( a,b \) and mutually independent sequences of identically distributed random variables \( (U_k)_k, (V_k)_k\) that are independent of \( X_0 \).
That is \( \mathbb{P}[Y_n \in A|Y_1 = y_{1}, \dots, Y_{n-1} = y_{n-1}, X_1 = x_1, \dots, X_n = x_n] = \mathbb{P}[Y_n \in A|X_n = x_n] \).
In order to model data and for the purpose of statistical inference, the model is considered in the filtration generated by 
\( (Y_k)_{k \geq 0} \). This means that the Markov process \((X_k)_{k \geq 0}\) is hidden.

MCMC has various applications in this context:
* Computing the posterior distribution of \(X_0, X_1, ..., X_n\) given observations \(Y_0 = y_0, ..., Y_n = y_n\).
* Other tasks: If the model is governed by unknown parameters, the parameter estimation may involve MCMC. 

For a concrete example, suppose you want to find the temperature \( \{X_n\}_{n \in S}, X_n \in \mathbb{R}\) in London, for each day during the summer of 1500,
but you only have information about the amount of water that was imported by the city each day \( \{Y_n\}_{n \in S}, Y_i \in \{1, 2, 3\} \), 
which is grouped into low, middle and a high amount. Then we have access to \( \mathbb{P}[Y_n = k|X_n = c] = \pi(Y_n|X_n)\) 
(simply by measuring the frequence of water \( k \)-amount water imported during the summer)
and an estimate for \( \mathbb{P}[X_n = c|Y_n = k] = p(X_n|Y_n) \) (the probability that it is \( c \)-degrees given that that \( k \)-amount water
was imported, so it should be an increasing function in \( k \)). 

Note that this places the restriction that weather may only depend on the previous day, instead of say the whole week.


Explain the Metropolis-Hastings algorithm in detail. How is it related to the sought
distribution?

For consistency reasons \( \Pi^{*} \) is referred here as \( \pi \) and \( q(x,y) \) as \( Q(y|x) \),
deviating slightly from the lecture.

Some prerequisites are required:
A Markov chain is said to be reversible if there exists a probability distribution \( \pi \) on \( \Omega \) such
that \( p(x | y)\pi(y) = p(y | x)\pi(x) \)
Further \( \pi \) is then the stationary/equilibrium/invariant distribution of a Markov chain
since
\( \sum_{x} p(y|x) \pi(x) = \sum_x p(x|y)\pi(y) = \pi(y) \).

Suppose we want to sample from a distribution \( \pi(x) := \frac{\tilde{\pi}(x)}{Z_p} \). To do this
we first construct a (reversible) Markov chain as follows. Let \( X_t = x \) be the current state. We then perform the
following two steps repeatedly:
1. Generate \( Y \sim Q(\cdot | x) \) for some Markov transition matrix \( Q \).
Let y be the generated value.
2. Set \( X_{t+1} = y \) with probability \( α(y | x) := 1 \wedge \frac{\tilde{\pi}(y)}{\tilde{\pi}(x)} \frac{Q(x|y)}{Q(y|x)} \)
Otherwise set \( X_{t+1} = x \).

The fraction in \( \alpha(y|x) \) is also referred to as Metropolis-Hastings ratio
and \( Q \) the proposal kernel.
The transition probabilities created by the algorithm (so not \( Q \)) are
\( p(y|x) = \alpha(y|x)Q(y|x) \) if \( y \neq x \) and \( p(y|x) = 1 - \sum_{y \neq x} \alpha(y|x)Q(y|x) \).

For the right choice of \( Q \) the resulting Markov chain is reversible 
and if \( \tilde{\pi}(x) > 0, Q(x|y) > 0, p(x|x) > 0  \) for some \( x \in \Omega \)
it is also irreducible and ergodic, so the ergodic theorem applies.
Then the stationary distribution is precisely \( \pi(x) = \tilde{\pi }\frac{x}{Z_p} \).

Note that \( Z_p \) is not required for the algorithm.

The rate of convergence to the target distribution is, however not easy to determine
concepts like exponential ergodicity are theoretically available but often not applicable in practice.

Proof: 
\( \pi \) is proporational by the constant \( Z_p \) to \( \tilde{\pi} \), since we can divide such terms out of the 
reversible Markov chain equation, it suffices to show that \( \tilde{\pi} \) is the distribuition which makes \( p(y|x) \) reversible.

\( p(y|x)\tilde{\pi}(x) = \alpha(y|X)Q(y|X)\tilde{\pi}(x) = (\frac{\tilde{\pi}(y)}{\tilde{\pi}(x)} \frac{Q(x|y)}{Q(y|x)} \wedge 1)Q(y|x)\tilde{\pi}(x) \)
\( =  Q(x|y)\tilde{\pi}(y) \wedge Q(y|x)\tilde{\pi}(x) = (1 \wedge \frac{\tilde{\pi}(x)}{\tilde{\pi}(y)} \frac{Q(y|x)}{Q(x|y)})Q(x|y)\tilde{\pi}(y) \)
\( \alpha(x|y)Q(x|y)\tilde{\pi}(y) = p(x|y)\pi(y) \)

How is a burn-in period used in the Metropolis-Hastings algorithm, and why?


The Metropolis-Hastings Algorithm - Burn-in Period

Convergence is an asymptotic property that characterizes the behavior of the samples, if their number approaches \(\infty\).
Typically, the Metropolis-Hastings chain \((\xi_{k})\) is started with some initial distribution and is only asymptotically stationary:

(The invariant distribution is the sought object that is asymptotically sampled with the Markov chain, but the invariant distribution is typically not available as a tractable initial distribution whenever MCMC needs to be used.)

For finite and not too large \(k\), the distribution of \(\xi_{k}\) might still deviate significantly from the invariant distribution.

Choose a burn-in period \(m \in \mathbb{N}\) to mitigate the impact of the wrong initial distribution.
Estimators are now redefined as follows:
\[
\pi(x) = \frac{1}{n-m} \sum_{k=m+1}^{n} 1_{\{\xi_{k}=x\}}
\]
\[
\hat{f}_{n} = \frac{1}{n-m} \sum_{k=m+1}^{n} f(\xi_{k})
\]

The burn-in period \(m\) is typically chosen as 10\% of the total run length in typical applications.


Define Exponential Ergodicity and relate it to MCMC

Exponential Ergodicity / Geometric ergodicity

Consider an arbitrary Markov chain \((\xi_{k})_{k \geq 0}\).

We consider the distance of one-dimensional marginal distributions from the invariant distribution:

\[ d_{k}(y) := \sum_{x \in E} |\mathbb{P}_{y}(\xi_{k} = x) - \pi(x)| \]

The Markov chain is exponentially ergodic, if 
\[ \exists \rho < 1, h: E \rightarrow (0, \infty): d_{k}(y) \leq h(y) \cdot \rho^{k} \]

The following is useful in the context of MCMC
Irreducible, aperiodic Markov chains on a finite state space are automatically exponentially ergodic.
The constant \(\rho\) is given by the absolute value of the eigenvalue of the one-step transition matrix having second-largest modulus.
In problems that require sophisticated simulation techniques like MCMC, analytical bounds on the second eigenvalue are typically too difficult to compute or too loose to be useful. 

Concrete implementations of the Metropolis-Hastings algorithm rely on different
proposal kernels. Discuss examples.

The Independence Sampler

The proposal is independent of the current state \(x\), i.e., 
\[ Q(y|x) = q(y),\ \forall x, y \in E \]
This special choice leads to a Metropolis-Hastings ratio

\[ \frac{\pi(y)}{\pi(x)} \frac{Q(x)}{Q(y)} = \frac{\omega(y)}{\omega(x)}\] 
with 
\[ \omega(x) = \frac{\pi(x)}{q(x)} \]


We define constants
\( C := \sum_{x \in E} \pi(x)\)
\(A := \sup_{x \in E} \omega(x) \)

The following conditions are equivalent:

1. The Metropolis-Hastings chain run by independent sampling is exponentially ergodic.
2. \( A < \infty \).

In this case,
\( d_k(y) \leq \left( 1 - \frac{C}{A} \right)^k \)

One can show that \( C \leq A \). The upper bound for the speed of convergence is the better, the closer \(A\) to \(C\).
It is not ideal, if \(Q(x)\) becomes much smaller than \(\pi(x)\) for some \(x\). For example, this is the case 
if \(E = \mathbb{Z}\) and \(Q\) is lighter-tailed than \(\pi\). 



Random-Walk proposal

We consider as an example \(E = \mathbb{Z}^d\).
The main idea is that the proposal \(Q(y|X)\) does not directly depend on the current location \(x\), but only on the difference 
between the future state \(y\) and the current state \(x\):

\( Q(y|x) = f(y - x) \) with \(\sum_{z \in \mathbb{Z}^d} f(z) = 1, f > 0\)

Let \(E = \mathbb{Z}\) and \(Q(y|x) = f(y - x)\).

Suppose that the function \(f\) is symmetric, i.e., \(f(z) = f(-z)\).
In addition, we make the following assumptions:
1. \( \exists \alpha > 0, x_0 > 0: \forall |y| \geq |x| \geq x_0 \):
2. \(\frac{\tilde{\pi(y)}}{\tilde{\pi}(x)} \leq e^{-\alpha(|y| - |x|)}\)
(This assumption implies that the tail of \(\tilde{\pi}\) decays at least exponentially fast.)
3. \( \forall x_1 > 0: \inf_{|x| \leq x_1} \tilde{\pi}(x) > 0 \).

Then the random-walk Metropolis-Hastings algorithm is exponentially ergodic.
